\documentclass[10pt,a4paper]{amsart}
\usepackage[margin=19mm]{geometry}
\usepackage{amsmath,amssymb,mathtools}
\usepackage[scr=rsfs,cal=euler]{mathalfa}
\usepackage{xfrac}
\usepackage[unicode,hidelinks,bookmarks=false]{hyperref}
\newcommand{\ie}{i.e.\ }
\newcommand{\cf}{cf.\ }
\newcommand{\resp}{respectively\ }
\newcommand{\Subsection}{\subsection}
\providecommand{\nobreakdash}{\nobreak\hbox{-}}
\setlength{\emergencystretch}{2em}
\multlinegap0pt
\usepackage{amssymb}
\usepackage{bm}
\newtheorem{lemma}{Lemma}
\newcommand{\R}{\mathbb R}
\newcommand{\la}{\lambda}
\newcommand{\ls}{\lesssim}
\title{Restriction estimates for toral eigenfunctions and lattice points in spherical regions}
\author{Cheng Zhang, Zhifei Zhu}
\date{Source: arXiv:2606.08650v1 (CC BY 4.0)}
\begin{document}
\maketitle

\noindent\emph{Source and licence.} Restriction estimates for toral eigenfunctions and lattice points in spherical regions. Version arXiv:2606.08650v1 (CC BY 4.0), distributed by arXiv under the Creative Commons Attribution 4.0 International licence, http://creativecommons.org/licenses/by/4.0/, as recorded in the arXiv licence metadata for this exact version and shown on the abstract page https://arxiv.org/abs/2606.08650v1. Full licence notice: \texttt{licenses/arxiv-2606.08650-CC-BY-4.0.txt}.\par\smallskip

\noindent\textit{Editorial scope.} Counted unit: the article's lemma lem:shell (source0.tex lines 920--940). The article's complete original proof of it is delivered with it (source0.tex lines 943--1063, one \texttt{proof} environment of 121 lines). No mathematics is written, repaired or re-derived: the statement and the proof are the article's own lines, in the article's own notation, and no retained line is substituted or re-flowed.  No further source-local context is needed: every cross-reference of every retained line resolves inside the two delivered spans.  One declared adaptation: exactly 1 ancillary strings inside the retained spans are omitted (image-inclusion commands and, where the source repeats a label, the repeated label definitions) and no other line differs from the source; the figure environments, with their placement, captions and labels, are kept, and the package records the omitted strings in fidelity receipts bound to the affected bodies.  No retained line contains a citation, so the wrapper carries no bibliography and no bibliography entry.  The retained lines contain the article's own diagram code, which the wrapper typesets with the standard tikz packages; no image file is delivered and the package declares no asset dependency.  Editorial formatting only, in addition: an AMS \texttt{amsart} preamble that restates the article's own declarations and macros used by the retained lines, namely the environments \texttt{lemma} on separate counters, together with the standard packages those lines need.  The retained environments print in this document as Lemma 1 in order of appearance, whereas the article prints the same items with the numbering of its own full document.  The complete native source of the article as distributed in the arXiv e-print for this exact version is bundled unchanged as \texttt{sources/202347/source0.tex}.
	\begin{lemma}\label{lem:shell}
		Let \(\R^d=V\oplus W\), with \(m=\dim W\ge2\). Suppose
		\[
		\Omega\subset\{x\in\la S^{d-1}:|\pi_Vx-p|\le C\}
		\]
		for some \(p\in V\).  Let \(R=(\max(\la^2-|p|^2,0))^{1/2}\).
		Then one of the following holds.
		\begin{enumerate}
			\item[(i)] If \(R\ls\la^{1/2}\), then
			\(\operatorname{diam}(\pi_W\Omega)\ls\la^{1/2}\).
			\item[(ii)] If \(R\gg \la^{1/2}\), then \(\Omega\) is covered by
			\(O(\la^{\frac{m-1}2})\) many sets \(\Omega_\alpha\), each satisfying
			\[
			\operatorname{diam}(\pi_{V\oplus\R\omega_\alpha}\Omega_\alpha)\ls1,
			\qquad
			\operatorname{diam}(\pi_{W\cap\omega_\alpha^\perp}\Omega_\alpha)
			\ls\la^{1/2}
			\]
			for some unit vector \(\omega_\alpha\in W\).
		\end{enumerate}
	\end{lemma}

	\begin{proof}
		Write every point $x\in\Omega$ as
		\[
		x=p+v+w,
		\qquad v\in V,
		\qquad w\in W,
		\qquad |v|\le C.
		\]
		Since $V\perp W$ and $|x|=\lambda$, we have
		\begin{equation}\label{eq:shell-basic-equation}
			|w|^2=\lambda^2-|p+v|^2.
		\end{equation}
		Therefore,
		\[
		\bigl||p+v|^2-|p|^2\bigr|
		\le 2|p|\,|v|+|v|^2
		\ls_C \lambda.
		\]
		By definition,
		\[
		R^2=\max(\lambda^2-|p|^2,0),
		\]
		this implies the uniform shell estimate
		\begin{equation}\label{eq:shell-square-thickness}
			\bigl||w|^2-R^2\bigr|\ls\lambda.
		\end{equation}
		Indeed, if $|p|\le\lambda$, then
		$R^2=\lambda^2-|p|^2$ and the conclusion follows directly from
		\eqref{eq:shell-basic-equation}.  If $|p|>\lambda$, then $R=0$ and the same
		bound follows because $\lambda^2-|p|^2<0$ while
		$|w|^2=\lambda^2-|p+v|^2\ge0$ differs from
		$\lambda^2-|p|^2$ by $O(\lambda)$.
		
		If $R\ls\lambda^{1/2}$, then \eqref{eq:shell-square-thickness} gives
		$|w|^2\ls\lambda$, hence $|w|\ls\lambda^{1/2}$ for all $x\in\Omega$.
		Thus,
		\[
		\operatorname{diam}(\pi_W\Omega)
		\le 2\sup_{x\in\Omega}|\pi_Wx|
		\ls\lambda^{1/2}.
		\]
		This proves (i).
		
		We now assume $R\gg\lambda^{1/2}$.  From
		\eqref{eq:shell-square-thickness},
		\[
		\bigl||w|-R\bigr|
		=\frac{\bigl||w|^2-R^2\bigr|}{|w|+R}
		\ls \frac{\lambda}{R}.
		\]
		Thus, $\pi_W\Omega$ lies in a radial shell around the sphere
		$\{w\in W:|w|=R\}$ of thickness $O(1+\lambda/R)$.
		
		Put $q=m-1=\dim W-1$.  Choose a maximal $c\lambda^{1/2}$-separated set on
		the sphere $\{w\in W:|w|=R\}$, where $c>0$ is a small fixed constant. The
		corresponding spherical caps of tangential radius $O(\lambda^{1/2})$ cover
		the sphere, and the number of such caps is
		$	O((\frac{R}{\lambda^{1/2}})^q),$
		because $R\gg\lambda^{1/2}$.  Let $R\omega$ be the center of one cap
		with $\omega\in W$ and $|\omega|=1$.
		
		Consider the portion of the shell that lies over this cap.  For such a point
		$w$, write
		\[
		w=(\omega\cdot w)\omega+w_\perp,
		\qquad w_\perp\in W\cap\omega^\perp.
		\]
		The angular localization gives
		\[
		|w_\perp|\ls\lambda^{1/2}.
		\]
		We will then estimate the range of the scalar coordinate $\omega\cdot w$.  Since
		$w$ is in a small cap, $\omega\cdot w$ is comparable to $R$, and
		\[
		|w|-\omega\cdot w
		=\frac{|w_\perp|^2}{|w|+\omega\cdot w}
		\ls \frac{\lambda}{R}.
		\]
		The radial variable $|w|$ itself varies over an interval of length
		$O(1+\lambda/R)$.  Hence $\omega\cdot w$ also varies over an interval of
		length $O(1+\lambda/R)$
		on this cap.  We partition that interval into subintervals of bounded
		length.  This introduces $O(1+\lambda/R)$ many subpieces for each angular cap.
		
		On each resulting piece $\Omega_\alpha$, the $V$-projection has bounded
		diameter as $|\pi_Vx-p|\le C$, and the $\omega$-coordinate has bounded
		diameter by construction. Therefore,
		\[
		\operatorname{diam}(\pi_{V\oplus\R\omega}\Omega_\alpha)\ls1.
		\]
		At the same time the transverse component in $W\cap\omega^\perp$ has
		diameter $O(\lambda^{1/2})$, so
		\[
		\operatorname{diam}(\pi_{W\cap\omega^\perp}\Omega_\alpha)
		\ls\lambda^{1/2}.
		\]
		See Figure \ref{fig1}.
		\begin{figure}[htbp]
			\centering
			
			\caption{Lemma \ref{lem:shell} (ii)}
			\label{fig1}
		\end{figure}
		
		It remains only to count the pieces.  Since $R\le\lambda+C\ls\lambda$,
		we have
		\[
		\Big(\frac{R}{\lambda^{1/2}}\Big)^q
		\Big(1+\frac{\lambda}{R}\Big)
		\ls
		\lambda^{q/2}.
		\]
		Indeed, the first term is at most $\lambda^{q/2}$, and the second is
		\[
		\Big(\frac{R}{\lambda^{1/2}}\Big)^q\frac{\lambda}{R}
		=R^{q-1}\lambda^{1-q/2}
		\ls\lambda^{q/2}.
		\]
		Thus $\Omega$ is covered by $O(\lambda^{q/2})=O(\lambda^{(m-1)/2})$
		pieces with the asserted projection bounds.  This proves (ii).
	\end{proof}

\end{document}
